% 原创算子图；两种 AD 模式共用节点位置，便于直接比较传播方向。
\begin{tikzpicture}[x=1mm,y=1mm,font=\figurefont\footnotesize,text=NNInk,
  ad input/.style={circle,draw=NNInput,fill=NNInput!18,line width=1.1pt,
    minimum size=11mm,inner sep=0pt},
  ad operator/.style={circle,draw=NNHidden,fill=NNHidden!18,line width=1.1pt,
    minimum size=11mm,inner sep=0pt},
  ad forward/.style={-{Stealth[length=1.8mm,width=1.3mm]},draw=NNInput,
    line width=1.2pt,rounded corners=2mm},
  ad derivative/.style={text=NNGradient,align=center,text width=30mm,
    inner sep=0pt}]
  \node[ad input] (u) at (12,14) {$u=2$};
  \node[ad input] (v) at (12,-14) {$v=3$};
  \node[ad derivative] at (12,3) {$\dot u=1$};
  \node[ad derivative] at (12,-25) {$\dot v=0$};
  \node[ad operator] (mul) at (52,0) {$\times$};
  \node[ad operator] (add) at (84,0) {$+$};
  \node[ad operator] (square) at (116,0) {$(\cdot)^2$};
  \node[ad operator,draw=NNOutput,fill=NNOutput!18] (half) at (148,0)
    {$\times\frac12$};
  \draw[ad forward] (u.east)--(mul.north west);
  \draw[ad forward] (v.east)--(mul.south west);
  \draw[ad forward] (u.north)--(12,28)--(84,28)--(add.north);
  \node[above,text=NNInput] at (49,28) {$u=2,\quad\dot u=1$};
  \draw[ad forward] (mul)--(add);
  \draw[ad forward] (add)--(square);
  \draw[ad forward] (square)--(half);
  \foreach \x/\value in {52/r=6,84/s=8,116/t=64,148/J=32}{
    \node[text=NNInput] at (\x,-12) {$\value$};
    \draw[NNLine,line width=.45pt] (\x-12,-20)--(\x+12,-20);
  }
  \node[ad derivative] at (52,-30) {$\begin{aligned}\dot r&=3\cdot1+2\cdot0\\&=3\end{aligned}$};
  \node[ad derivative] at (84,-30) {$\begin{aligned}\dot s&=3+1\\&=4\end{aligned}$};
  \node[ad derivative] at (116,-30) {$\begin{aligned}\dot t&=2\cdot8\cdot4\\&=64\end{aligned}$};
  \node[ad derivative] at (148,-30) {$\begin{aligned}\dot J&=\tfrac12\cdot64\\&=32\end{aligned}$};
  \draw[ad forward] (22,-47)--(32,-47);
  \node[anchor=west] at (35,-47) {计算依赖};
  \node[anchor=west,text=NNInput] at (82,-47) {$r=6$};
  \node[anchor=west] at (97,-47) {值};
  \node[anchor=west,text=NNGradient] at (111,-47) {$\dot r=3$};
  \node[anchor=west] at (127,-47) {方向导数};
\end{tikzpicture}
